%%%%%%%%%%%%%%%%%%%%%%%%%%%%%%%%%%%%%%%%%%%%%%%%%%%%%%%%%%%%%%%%%%%%%%%%%%%
% Lehreinheit 12: Präsentation und Glaubwürdigkeit
% Quelle: WuZ-Buch, Kap. 15 (LE 12), Kap. 6 Persona/CREDS (sec:creds,
%         Goldener Ablauf, Nocebo der Lehre, sec:lernplacebo),
%         Kap. 9 Skript-Gerüst (ch:skript), Kap. 3d Dr.-Fox-Effekt,
%         Anhang C Magiculum (Kontext und Alignment)
% Staunmoment: CRED-Strategien des Zauberkünstlers  |  Kernbegriff: Credibility Cues
%%%%%%%%%%%%%%%%%%%%%%%%%%%%%%%%%%%%%%%%%%%%%%%%%%%%%%%%%%%%%%%%%%%%%%%%%%%
\nsl{Wer Forschungsergebnisse vorträgt, betreibt Überzeugungsarbeit, ob er
will oder nicht. Diese Einheit beginnt mit zwei kurzen Vortragsausschnitten:
Der eine wirkt überzeugend, der andere ist fachlich stärker. Die Gruppe
sammelt, was den Eindruck macht, und trennt dabei Auftreten von Substanz.
Das Werkzeug dafür liefert das Buch: Credibility Enhancing Displays
(CREDS), also Handlungen und Signale, die Glaubwürdigkeit erzeugen. Die
Einheit verwendet sie als Analyseinstrument und warnt zugleich vor ihnen,
denn dieselben Signale erzeugen auch Scheinsicherheit. Danach geht es um das
Handwerk: Dramaturgie eines Kurzvortrags, Foliengestaltung, Transparenz
über Grenzen, Umgang mit Fragen, Poster. In der zweiten Hälfte halten die
Studierenden ihre 5-Minuten-Präsentation zum Semesterprojekt.
\newpage}
\os{\begin{description}
\item[{\rot\bf Heute}]~\\[-\lblskip]
\bi
  \ite zwei Vorträge: Wem glauben Sie -- und warum?
  \ite Credibility Cues: Werkzeug und Warnung
  \ite Dramaturgie eines Kurzvortrags
  \ite Folien, die das Arbeitsgedächtnis schonen
  \ite Limitationen zeigen, Fragen beantworten
  \ite Ihre 5-Minuten-Präsentationen
\ei
\end{description}
\newpage}

\begin{vorbereitung}[CRED-Strategien des Zauberkünstlers]
\textbf{Eine Woche vorher}
\begin{checkliste}
  \item \textbf{Zwei Vortragsausschnitte auswählen} (je 2--3 Minuten).
    Kriterien: (a) Ausschnitt~A wirkt überzeugend (flüssig, sicher,
    anschaulich), ist aber fachlich schwach: Er verallgemeinert, nennt keine
    Grenzen oder vertritt einen Befund, der sich später nicht bestätigt hat.
    (b) Ausschnitt~B ist fachlich stärker (Daten, Einschränkungen,
    saubere Begriffe), wirkt aber weniger gewinnend. (c) Beide behandeln
    möglichst dasselbe oder ein verwandtes Thema, damit sich nur
    Auftreten und Substanz unterscheiden. (d) Sie selbst kennen den
    fachlichen Stand zu beiden Ausschnitten und haben Belege dafür
    bereit. (e) Niemand wird vorgeführt: Personen in öffentlicher,
    beruflicher Rolle, Analyse des Vortrags, nicht der Person.
  \item \textbf{Alternative ohne Rechtefragen:} beide Ausschnitte selbst
    aufnehmen, mit identischem Inhaltskern -- einmal mit starker Performance
    und überzogenen Aussagen, einmal mit vorsichtigen Aussagen und
    Limitationen, aber flachem Vortrag. Das ist die Logik des
    Dr.-Fox-Experiments \parencite{naftulin1973drfox} und zugleich ein
    sauberes Ein-Stimulus-Design (LE~6).
  \item \textbf{Rechtliches} (keine Rechtsberatung; im Zweifel die
    Rechtsabteilung der Hochschule fragen): Das Zitatrecht (§\,51 UrhG)
    deckt kurze Ausschnitte, wenn sie belegt und inhaltlich analysiert
    werden und die Quelle genannt wird. §\,60a UrhG erlaubt für
    Lehrveranstaltungen die Nutzung von Teilen veröffentlichter Werke
    (bis 15\,\%) im Kreis der Teilnehmenden. Frei zugängliche Videos am
    besten direkt von der Originalplattform abspielen statt herunterladen
    und weiterverbreiten. Bei Videos unter Creative-Commons-Lizenz die
    Lizenzbedingungen prüfen (z.\,B.\ Namensnennung, keine Bearbeitung);
    Ausschnitte nicht in eine öffentlich zugängliche Folienfassung
    einbetten, sondern dort nur verlinken.
  \item Studierende erinnern: Präsentation höchstens 5 Minuten, Folien
    als PDF bis zum Vortag hochladen (ein Laptop für alle spart
    Wechselzeit).
  \item Zeitplan für die Präsentationsrunde festlegen (siehe
    Lehrhinweis in Abschnitt~\ref{sec:le12-praes}) und die Reihenfolge
    vorab bekanntgeben.
\end{checkliste}

\textbf{Handout und Übungsblatt}
\begin{checkliste}
  \item Handout \enquote{Präsentation und Glaubwürdigkeit} und Übungsblatt
    LE~12 je einmal pro Person kopieren.
  \item Feedbackbogen (Aufgabe~2) in ausreichender Zahl: ein Bogen pro
    Person reicht für fünf Vorträge; bei mehr Vorträgen zwei Bögen.
  \item Bewertungsraster (Aufgabe~3) für die Lehrperson: ein Exemplar pro
    Vortrag.
  \item Die drei Beispielfolien für die Folienkritik (Aufgabe~4) zusätzlich
    als Folien bereitlegen, damit sie groß gezeigt werden können.
\end{checkliste}

\textbf{Am Tag selbst}
\begin{checkliste}
  \item Beide Videos vorher am Hörsaalrechner testen (Ton!); Plan B: lokale
    Kopie für den Notfall nur, wenn die Rechte das erlauben.
  \item Gut sichtbare Uhr oder Timer für die Vortragenden; Signal bei
    4:00 Minuten, Schluss bei 5:00.
  \item Tafel für den Vergleich vorbereiten: zwei Spalten
    \emph{wirkt überzeugend, weil \ldots} und \emph{ist belastbar, weil
    \ldots}.
\end{checkliste}

\textbf{Risiken und Plan B.} Die Gruppe hält oft den Vortrag für fachlich
stärker, der besser wirkt; das ist der erwünschte Befund und sollte nicht
vorzeitig korrigiert werden. Erst die Bewertungen einsammeln, dann die
Auflösung. Wenn Studierende aufgeregt sind: Die Präsentation ist ein
Zwischenbericht, offene Punkte sind ausdrücklich erlaubt. Den
Nocebo-Satz \enquote{Das ist für viele hier schwer} vermeiden.
\end{vorbereitung}

\begin{ablauf}
\begin{description}
\item[0--15 · Staunmoment/Analyse.] Zwei Vortragsausschnitte im Vergleich.
  Alle bewerten jeden Ausschnitt auf zwei Skalen (\emph{überzeugend},
  \emph{fachlich belastbar}) und notieren, woran sie das festmachen. Danach
  Auflösung und Sammlung an der Tafel: Was ist Performance, was Substanz?
\item[15--40 · Theorie-Input.] CREDS als Analysewerkzeug und ihre
  Scheinsicherheit; Vortragsdramaturgie und Goldener Ablauf;
  Foliengestaltung (Signal-Rausch, Dual Coding, Arbeitsgedächtnis);
  Transparenz über Limitationen und das Nocebo der Lehre; Umgang mit
  Fragen; Poster.
\item[40--75 · Kurzpräsentationen.] 5-Minuten-Vorträge zum
  Semesterprojekt mit Feedback nach dem Raster \emph{klar -- überzeugt --
  offen} (Aufgabe~2).
\item[75--90 · Folienkritik.] Drei Beispielfolien werden nach Klarheit,
  Arbeitsgedächtnis-Last und Botschaftskonsistenz überarbeitet
  (Aufgabe~4).
\end{description}
\end{ablauf}

%==============================================================================
\section{Staunmoment: Zwei Vorträge}
%==============================================================================
\begin{staunbox}{Wem glauben Sie?}
Sie sehen zwei kurze Ausschnitte aus Vorträgen.

\medskip
\emph{Bewerten Sie jeden Ausschnitt sofort danach auf zwei Skalen von 1 bis 5:
Wie überzeugend wirkt er? Wie belastbar ist das Gesagte fachlich? Notieren
Sie zu jeder Zahl einen Grund.}
\end{staunbox}
\os{\newpage}

\begin{description}
\item[{\rot\bf Was ist gerade passiert?}]~\\[-\lblskip]
\bi
  \ite Welcher Ausschnitt hat Sie mehr überzeugt? Woran lag das?
  \ite Welcher ist fachlich belastbarer? Woher wissen Sie das?
  \ite Liegen Ihre beiden Bewertungen nah beieinander -- oder weit
    auseinander?
\ei
\os{\newpage}

\item[{\rot\bf Performance oder Substanz?}]~\\[-\lblskip]
\begin{center}\nsl{\small}
\begin{tabular}{p{0.42\textwidth}p{0.42\textwidth}}
\toprule
\textbf{wirkt überzeugend, weil \ldots} & \textbf{ist belastbar, weil \ldots} \\
\midrule
sicherer Blick, feste Stimme & Daten werden gezeigt \\
flüssige Sprache, keine Pausen & Quelle und Methode sind genannt \\
anschauliche Geschichte & Grenzen werden benannt \\
Gewissheit (\enquote{immer}, \enquote{nachweislich}) & Gewissheit passt zu den Daten \\
\bottomrule
\end{tabular}
\end{center}
\bi
  \ite die linke Spalte sagt etwas über den Vortrag, die rechte über die
    Aussage
\ei
\os{\newpage}
\end{description}

\begin{lehrhinweis}
Die Bewertungen per Handzeichen oder Zettel einsammeln, bevor Sie etwas
auflösen, und den Mittelwert beider Skalen für beide Ausschnitte an die
Tafel schreiben. Der typische Befund -- A wird auf beiden Skalen besser
bewertet, obwohl B fachlich stärker ist -- ist der Dr.-Fox-Effekt im
Kleinen. Die Tabelle oben ist eine erwartbare Sammlung; die Gruppe wird
eigene Punkte finden. Danach den Satz aus dem Buch anschreiben:
\enquote{Wir verwechseln Auftreten mit Wahrheit und Bühne mit Beweis.}
\end{lehrhinweis}

%==============================================================================
\section{Präsentieren heißt überzeugen}
%==============================================================================
\nsl{Eine wissenschaftliche Präsentation ist ein kommunikativer Akt mit
Publikum, Kontext und Rolle. Markus Pfeil beschreibt für die Zauberkunst,
wie stark die Wirkung einer Vorführung davon abhängt, ob die
\emph{Persona} des Vortragenden zum \emph{Kontext} passt, in dem das
Publikum ihn sieht \parencite{pfeil2019}. Für Vortragende in der
Wissenschaft gilt dasselbe: Eine Studentin im Seminar, eine Professorin auf
einer Tagung und ein Experte in einer Talkshow werden mit
unterschiedlichen Erwartungen gehört. Frieder Vogelmann beschreibt, wie
Wissen nicht allein durch seinen Gehalt wirksam wird, sondern durch
institutionelle Einbettung und Inszenierung
\parencite{vogelmann2020wirksamkeit}. Der Zauberkünstler erzeugt im
Extremfall eine Wahrheit allein durch Inszenierung. Die Wissenschaft muss
sich von dieser Form abheben, indem sie ihr Verfahren offenlegt.}
\begin{description}
\item[{\rot\bf Der Vortrag als Bühne}]~\\[-\lblskip]
\bi
  \ite \textbf{Kontext}: Wo spreche ich, was weiß das Publikum über mich?
  \ite \textbf{Persona}: Welche Rolle nehme ich ein -- und passt sie?
    \parencite{pfeil2019}
  \ite Wissen wird auch durch Inszenierung wirksam
    \parencite{vogelmann2020wirksamkeit}
  \ite Unterschied zur Bühne: Die Wissenschaft legt ihr Verfahren offen
\ei
\os{\newpage}

\item[{\rot\bf Der Dr.-Fox-Effekt}]~\\[-\lblskip]
\bi
  \ite 1973: ein Schauspieler als \enquote{Dr.\ Myron L.\ Fox}, Vortrag
    mitreißend, Inhalt bewusst sinnlos \parencite{naftulin1973drfox}
  \ite das Fachpublikum bewertete ihn glänzend
  \ite verschärfte Replikation: Charisma hebt Zufriedenheit und
    \emph{gefühlte} Lernleistung, nicht das Gelernte
    \parencite{peerbabad2014drfox}
  \ite Folgerung: Ihr Publikum ist gegen gute Performance nicht immun --
    Sie auch nicht
\ei
\os{\newpage}
\end{description}

%==============================================================================
\section{CREDS: Werkzeug und Warnung}
%==============================================================================
\nsl{Der Anthropologe Joseph Henrich hat beschrieben, warum Menschen nicht
einfach glauben, was jemand sagt, sondern darauf achten, was er
\emph{tut}: Worte sind billig, Handlungen mit Kosten oder Risiko sind
schwer zu fälschen \parencite{henrich2009}. Solche
\emph{Credibility Enhancing Displays} (CREDS) nutzt die Zauberkunst
systematisch, etwa durch erschwerende Bedingungen, das Untersuchenlassen
der Requisiten oder eine sichtbar geübte, flüssige Handhabung. Auch ein
Vortrag ist voller solcher Signale. Sie zu benennen hilft bei der Analyse
fremder und eigener Vorträge. Die Warnung gehört dazu: Ein Signal kann
Glaubwürdigkeit erzeugen, ohne dass die Aussage sie verdient. Schon die
leichte Lesbarkeit einer Aussage lässt sie wahrer erscheinen
\parencite{reber1999fluency}.}
\begin{description}
\item[{\rot\bf Was sind CREDS?}]~\\
\defn{credibilitycue}{Ein \textbf{Credibility Cue} ist ein Signal, an dem
ein Publikum die Glaubwürdigkeit einer Person oder Aussage abliest.
\textbf{CREDS} (Credibility Enhancing Displays) sind solche Signale in Form
von Handlungen, die Kosten oder Risiko tragen \parencite{henrich2009}.}
\bi
  \ite Worte sind billig, Handlungen nicht
\ei
\os{\newpage}

\item[{\rot\bf CREDS im Vortrag}]~\\[-\lblskip]
\begin{center}\nsl{\small}
\begin{tabular}{p{0.3\textwidth}p{0.28\textwidth}p{0.3\textwidth}}
\toprule
\textbf{Signal} & \textbf{wirkt als} & \textbf{belegt tatsächlich} \\
\midrule
flüssiger, freier Vortrag & Kompetenz & Übung \\
Fachbegriffe, Titel & Expertise & Zugehörigkeit \\
eigene Rohdaten gezeigt & Offenheit & prüfbares Material \\
Grenzen benannt & Ehrlichkeit & Kenntnis der Methode \\
\bottomrule
\end{tabular}
\end{center}
\bi
  \ite nur die unteren beiden sind \emph{kostspielig} -- sie machen
    angreifbar
\ei
\os{\newpage}

\item[{\rot\bf Die Warnung: Scheinsicherheit}]~\\[-\lblskip]
\bi
  \ite leicht Lesbares wird eher für wahr gehalten
    \parencite{reber1999fluency}
  \ite charismatischer Vortrag $\neq$ mehr Gelerntes
    \parencite{peerbabad2014drfox}
  \ite CREDS sagen, wie \emph{überzeugt} jemand ist -- nicht, ob er
    \emph{recht} hat
  \ite Prüffrage für jeden Cue: Was würde ich sehen, wenn die Aussage
    falsch wäre?
\ei
\os{\newpage}
\end{description}

\begin{merkbox}[Credibility Cues]
Glaubwürdigkeit entsteht im Publikum; Belastbarkeit liegt in der Aussage.
Gute Wissenschaftskommunikation sorgt dafür, dass beides
zusammenpasst -- und wer zuhört, prüft, ob es das tut.
\end{merkbox}
\os{\newpage}

%==============================================================================
\section{Dramaturgie eines Kurzvortrags}
%==============================================================================
\nsl{Das Buch beschreibt für den Einsatz eines Zauberkunststücks im
Unterricht einen \enquote{Goldenen Ablauf}: thematischer Einstieg, ein
glaubwürdigkeitsstiftendes Signal, der Effekt, eine Pause von etwa zehn
Sekunden, dann der Transfer zum Thema. Das Skript-Gerüst des Buchs
ordnet dieselben Schritte in vier Phasen: Framing, CRED, Vorführung,
Brücke. Ein wissenschaftlicher Kurzvortrag lässt sich nach demselben Muster
bauen. An die Stelle des Effekts tritt der zentrale Befund, an die Stelle
des Requisits die Methode. Die Pause nach dem Befund ist kein Zufall:
Nach dem Buch entsteht Lernbereitschaft gerade in dem Moment, in dem eine
offene Frage im Raum steht.}
\begin{description}
\item[{\rot\bf Vom Goldenen Ablauf zum Kurzvortrag}]~\\[-\lblskip]
\begin{center}\nsl{\small}
\begin{tabular}{p{0.38\textwidth}p{0.5\textwidth}}
\toprule
\textbf{Goldener Ablauf (Buch)} & \textbf{5-Minuten-Vortrag} \\
\midrule
thematischer Einstieg & Frage und warum sie zählt \\
das CRED & Methode: wie ich es geprüft habe \\
der Effekt & der zentrale Befund \\
die Aha-Pause & Pause, Befund wirken lassen \\
der Transfer & Bedeutung, Grenzen, nächste Frage \\
\bottomrule
\end{tabular}
\end{center}
\os{\newpage}

\item[{\rot\bf Fünf Minuten, eine Botschaft}]~\\[-\lblskip]
\begin{center}
\resizebox{0.95\textwidth}{!}{%
\begin{tikzpicture}[font=\sffamily]
  \foreach \x/\w/\t/\c in {0/1/Frage/rwuviolet!15,1/0.75/Lücke/rwuviolet!15,
      1.75/1/Methode/rwucyan40!50,2.75/1.25/Befunde/rwucyan40!80,
      4/0.5/Grenzen/rwuviolet!15,4.5/0.5/Botschaft/rwuviolet!35}{
    \fill[\c] (\x*3,0) rectangle ({(\x+\w)*3},1.1);
    \draw[rwuviolet] (\x*3,0) rectangle ({(\x+\w)*3},1.1);
    \node[font=\sffamily\footnotesize] at ({(\x+\w/2)*3},0.55) {\t};}
  \foreach \m in {0,1,2,3,4,5}{
    \draw (\m*3,0) -- (\m*3,-0.2) node[below,font=\sffamily\scriptsize]{\m:00};}
\end{tikzpicture}}
\end{center}
\bi
  \ite eine Botschaft, die man in einem Satz sagen kann
  \ite etwa eine Folie pro Minute; keine Folie ohne Aufgabe
  \ite den ersten und den letzten Satz auswendig können
\ei
\os{\newpage}

\item[{\rot\bf Der rote Faden}]~\\[-\lblskip]
\bi
  \ite jede Folie beantwortet eine Frage, die die vorige aufgeworfen hat
  \ite Überleitungen laut aussprechen: \enquote{Wie haben wir das
    gemessen?}
  \ite eine Pause nach dem Befund -- nicht sofort erklären
  \ite Ende = Botschaft, nicht \enquote{Danke für die Aufmerksamkeit}
\ei
\os{\newpage}
\end{description}

%==============================================================================
\section{Folien gestalten}
%==============================================================================
\nsl{Folien muss das Publikum gleichzeitig mit dem gesprochenen Wort
verarbeiten. Drei Befunde der
Lernpsychologie sind dafür maßgeblich. Erstens ist das Arbeitsgedächtnis
eng begrenzt; wenn Gruppierung und Wiederholung ausgeschlossen sind, hält
es nur etwa drei bis fünf Einheiten zugleich \parencite{cowan2001magical}.
Was es unnötig belastet, fehlt für das Verstehen
\parencite{sweller1988cognitive}. Zweitens werden sprachliche und bildhafte
Information in getrennten, aber verbundenen Systemen verarbeitet; was in
beiden kodiert ist, bleibt besser haften \parencite{paivio1986mental}.
Drittens hat Richard Mayer daraus Gestaltungsprinzipien abgeleitet, etwa
Unwesentliches wegzulassen, Wesentliches hervorzuheben und gesprochenen
Text nicht zusätzlich vollständig einzublenden \parencite{mayer2009multimedia}.
Edward Tufte hat für Diagramme das Verhältnis von Daten-Tinte zu
Dekoration als Gütemaß vorgeschlagen \parencite{tufte2001visual}, und
Michael Alley empfiehlt, die Folienüberschrift als vollständige Aussage zu
formulieren und darunter den Beleg zu zeigen \parencite{alley2013craft}.}
\begin{description}
\item[{\rot\bf Drei Grundlagen}]~\\[-\lblskip]
\bi
  \ite \textbf{Arbeitsgedächtnis}: etwa 3--5 Einheiten zugleich
    \parencite{cowan2001magical}; jede überflüssige Information kostet
    \parencite{sweller1988cognitive}
  \ite \textbf{Dual Coding}: Wort \emph{und} Bild werden getrennt verarbeitet
    und gemeinsam besser behalten \parencite{paivio1986mental}
  \ite \textbf{Signal-Rausch}: so viel Tinte wie möglich für Daten, so
    wenig wie möglich für Dekoration \parencite{tufte2001visual}
\ei
\os{\newpage}

\item[{\rot\bf Gestaltungsregeln}]~\\[-\lblskip]
\bi
  \ite \textbf{Kohärenz}: alles weg, was nicht zur Botschaft gehört
  \ite \textbf{Signalisierung}: das Wichtige hervorheben, nur das
  \ite \textbf{Redundanz}: nicht vorlesen, was auf der Folie steht
  \ite \textbf{Nähe}: Beschriftung an die Kurve, nicht in die Legende
  \ite \parencite{mayer2009multimedia}
\ei
\os{\newpage}

\item[{\rot\bf Überschrift als Aussage}]~\\[-\lblskip]
\begin{center}
\resizebox{0.9\textwidth}{!}{%
\begin{tikzpicture}[font=\sffamily]
  \draw[rwugray,thick] (0,0) rectangle (6,4);
  \node[anchor=north west,font=\sffamily\bfseries] at (0.2,3.85) {Ergebnisse};
  \foreach \y in {3.0,2.6,2.2,1.8,1.4,1.0,0.6}
    \draw[rwugray!70,line width=2pt] (0.4,\y) -- (5.4,\y);
  \node[below,font=\sffamily\small] at (3,-0.1) {Thema als Titel};
  \begin{scope}[xshift=7.5cm]
  \draw[rwucyan,thick] (0,0) rectangle (6,4);
  \node[anchor=north west,align=left,font=\sffamily\bfseries,text width=5.6cm] at (0.2,3.95)
    {Nach dem Staunmoment wurden mehr Begriffe erinnert};
  \fill[rwuviolet!40] (1.4,0.4) rectangle (2.4,0.9);
  \fill[rwucyan] (3.4,0.4) rectangle (4.4,1.4);
  \draw (1,0.4) -- (5,0.4);
  \node[font=\sffamily\scriptsize] at (1.9,0.2) {ohne};
  \node[font=\sffamily\scriptsize] at (3.9,0.2) {mit};
  \node[below,font=\sffamily\small] at (3,-0.1) {Aussage als Titel, Beleg darunter};
  \end{scope}
\end{tikzpicture}}
\end{center}
\bi
  \ite Assertion-Evidence-Prinzip \parencite{alley2013craft}; die Daten
    im Beispiel sind fiktiv
\ei
\os{\newpage}
\end{description}

%==============================================================================
\section{Transparenz über Limitationen}
%==============================================================================
\nsl{Eine wissenschaftliche Präsentation unterscheidet sich von einer
Zaubervorführung an einem Punkt grundsätzlich: Sie legt ihre Grenzen offen.
Das kostet scheinbar Überzeugungskraft. Tatsächlich ist das offene
Benennen von Einschränkungen selbst ein kostspieliges Signal im Sinne
Henrichs, denn es macht angreifbar, und gerade deshalb erhöht es die
Glaubwürdigkeit bei einem Publikum, das prüfen kann. Davon zu unterscheiden
ist die Selbstabwertung. Das Buch beschreibt das \enquote{Nocebo der
Lehre}: Negative Ankündigungen über die Lernfähigkeit von Studierenden
setzen Erwartungen, die sich auf Aufmerksamkeit und Arbeitsgedächtnis
auswirken können. Im Vortrag wirkt die entschuldigende Einleitung
(\enquote{Das ist jetzt vermutlich langweilig \ldots}) ähnlich. Grenzen
benennt man sachlich und an der richtigen Stelle, nämlich nach dem Befund.}
\begin{description}
\item[{\rot\bf Grenzen zeigen, ohne sich klein zu machen}]~\\[-\lblskip]
\begin{center}\nsl{\small}
\begin{tabular}{p{0.42\textwidth}p{0.46\textwidth}}
\toprule
\textbf{Nocebo (vermeiden)} & \textbf{Limitation (benennen)} \\
\midrule
\enquote{Das ist sicher langweilig.} & \enquote{Die Stichprobe umfasst 18 Personen aus einem Kurs.} \\
\enquote{Ich bin da nicht so sicher \ldots} & \enquote{Ob der Effekt anhält, haben wir nicht gemessen.} \\
\enquote{Das kriegen Sie eh nicht mit.} & \enquote{Eine Alternativerklärung wäre \ldots} \\
\bottomrule
\end{tabular}
\end{center}
\nsl{Die Zahl im Beispielsatz ist erfunden.}
\os{\newpage}

\item[{\rot\bf Warum Limitationen überzeugen}]~\\[-\lblskip]
\bi
  \ite sie sind \emph{kostspielig}: sie liefern dem Publikum Angriffspunkte
  \ite sie zeigen, dass Sie Ihre Methode kennen
  \ite sie halten die Frage offen -- das Publikum denkt weiter
  \ite Mara im Buch: Die leise Vortragende, die \enquote{wir wissen es
    noch nicht} sagte, behielt am Ende recht
\ei
\os{\newpage}
\end{description}

%==============================================================================
\section{Umgang mit Fragen}
%==============================================================================
\nsl{Fragen nach dem Vortrag sind organisierte Skepsis im Kleinen. Das Buch bietet für den Unterricht einen \enquote{Notfall-Satz}
an, wenn jemand ruft, er wisse, wie der Trick funktioniere: den Einwurf als
Hypothese würdigen und gemeinsam prüfen, ob er die Beobachtung vollständig
erklärt. Dieselbe Haltung trägt im wissenschaftlichen Vortrag. Hilfreich
ist es, Fragen zu sortieren: Verständnisfragen beantwortet man kurz,
Einwände nimmt man als Hypothese ernst, und Ko-Referate, die eigentlich
eigene Positionen vortragen, darf man freundlich begrenzen.}
\begin{description}
\item[{\rot\bf Vier Regeln für die Fragerunde}]~\\[-\lblskip]
\bi
  \ite Frage wiederholen oder umformulieren -- alle sollen sie hören
  \ite Einwand als Hypothese behandeln: \enquote{Prüfen wir, ob das die
    Daten erklärt.}
  \ite \enquote{Das weiß ich nicht} ist eine zulässige Antwort, mit
    \enquote{das wäre zu prüfen durch \ldots}
  \ite eine Backup-Folie mit Details für die wahrscheinlichste Frage
\ei
\os{\newpage}
\end{description}

%==============================================================================
\section{Poster und Laienpublikum}
%==============================================================================
\nsl{Das wissenschaftliche Poster ist ein eigenes Format. Es wird im
Vorbeigehen gelesen, meist in wenigen Sekunden entschieden, ob man stehen
bleibt, und dann im Gespräch erschlossen. Daraus folgen andere Regeln als
für Folien: eine Überschrift, die den Befund nennt; ein klarer Leseweg
(IMRaD in Spalten); große Abbildungen; wenig Fließtext. Der
Posterautor hält eine Zwei-Minuten-Fassung seines Vortrags bereit. Für ein
Laienpublikum gelten die Regeln dieser Einheit verschärft: Fachbegriffe
erklären oder weglassen, Unsicherheit trotzdem nicht verschweigen.}
\begin{description}
\item[{\rot\bf Poster}]~\\[-\lblskip]
\bi
  \ite Titel = Befund: aus 2\,m Entfernung lesbar
  \ite Leseweg in Spalten: Frage -- Methode -- Ergebnis -- Bedeutung
  \ite eine zentrale Abbildung, wenig Text
  \ite Zwei-Minuten-Fassung für das Gespräch bereithalten
\ei
\os{\newpage}

\item[{\rot\bf Für Laien sprechen}]~\\[-\lblskip]
\bi
  \ite ein Beispiel vor jeder Abstraktion
  \ite Fachbegriffe erklären oder ersetzen -- nicht weglassen, was sie
    bedeuten
  \ite Unsicherheit in Alltagssprache: \enquote{Das spricht dafür, ist
    aber noch nicht gesichert.}
  \ite Vorsicht vor dem Dr.-Fox-Effekt: Begeisterung ist kein Verstehen
\ei
\os{\newpage}
\end{description}

%==============================================================================
\section{Übung in der Stunde}
\label{sec:le12-praes}
%==============================================================================
\begin{description}
\item[{\rot\bf Ihre 5-Minuten-Präsentation}]~\\[-\lblskip]
\bi
  \ite Forschungsfrage, Methode, bisherige Befunde, offene Punkte
  \ite höchstens 5 Minuten; Signal bei 4:00
  \ite danach 2 Minuten Feedback: Was war \textbf{klar}? Was hat
    \textbf{überzeugt}? Was blieb \textbf{offen}?
  \ite alle anderen füllen den Feedbackbogen aus (Aufgabe~2)
\ei
\os{\newpage}

\item[{\rot\bf Folienkritik (15 Min.)}]~\\[-\lblskip]
\bi
  \ite drei Beispielfolien auf dem Übungsblatt (Aufgabe~4)
  \ite Kriterien: Klarheit, Last für das Arbeitsgedächtnis, konsistente
    Botschaft
  \ite für jede Folie: Was ist das Problem? Wie sähe die bessere Fassung
    aus?
\ei
\os{\newpage}
\end{description}

\begin{lehrhinweis}[Zeitplan der Präsentationsrunde]
Pro Vortrag sind realistisch 5 Minuten Vortrag, 2 Minuten mündliches
Feedback und eine halbe Minute Wechsel zu rechnen, also 7,5 Minuten. In
das 35-Minuten-Fenster passen damit \textbf{vier, knapp fünf Vorträge}.
Das reicht nur für sehr kleine Gruppen. Möglichkeiten:

\begin{center}\small
\begin{tabular}{p{0.2\textwidth}p{0.72\textwidth}}
\toprule
\textbf{Vorträge} & \textbf{Organisation} \\
\midrule
bis 5 & Plenum, wie im Ablauf vorgesehen. \\
6--7 & Folienkritik als Hausaufgabe auslagern; das Fenster wächst auf
  50 Minuten (40--90), das reicht für 6--7 Vorträge. \\
8--15 & Semesterprojekte in Zweier- oder Dreierteams präsentieren lassen
  (ein Vortrag pro Team), oder zwei parallele Gruppen in getrennten Räumen
  mit studentischer Moderation und Zeitnahme; die Lehrperson wechselt. \\
mehr als 15 & Zusätzlichen Präsentationstermin ansetzen oder parallele
  Gruppen; zur Bewertung Aufzeichnung nur mit schriftlicher Einwilligung
  der Vortragenden (vgl.\ LE~13). \\
\bottomrule
\end{tabular}
\end{center}

Da die Präsentation 15\,\% der Modulnote ausmacht, muss die Lehrperson
jeden Vortrag sehen oder hören können. Bei parallelen Gruppen ist das ohne
Aufzeichnung oder zweiten Termin nicht gewährleistet. Die Reihenfolge vorab
festlegen, alle Folien auf einem Rechner sammeln, einen Zeitnehmer bestimmen.
Das mündliche Feedback auf je einen Punkt pro Kategorie (klar, überzeugt,
offen) begrenzen; ausführlicher wird es auf den Bögen, die die Vortragenden
mitnehmen.
\end{lehrhinweis}

\begin{merkbox}
Ein guter wissenschaftlicher Vortrag überzeugt mit Mitteln, die er selbst
offenlegt: Frage, Methode, Befund und Grenzen. Was darüber hinaus
überzeugt, verdient einen zweiten Blick.
\end{merkbox}
\os{\newpage}

%==============================================================================
\section{Lernziele und Aufgaben}
%==============================================================================
\begin{description}
\item[{\rot\bf Nach dieser Einheit können Sie}]~\\[-\lblskip]
\bi
  \ite einen Forschungsbefund in fünf Minuten strukturiert präsentieren
  \ite Credibility Cues in Vorträgen erkennen und von Belegen
    unterscheiden
  \ite Folien nach Klarheit, Arbeitsgedächtnis-Last und Botschaft
    gestalten und kritisieren
  \ite Limitationen benennen, ohne sich abzuwerten
  \ite mit Fragen und Kritik sachlich umgehen
\ei
\os{\newpage}

\item[{\rot\bf Bis zur nächsten Woche}]~\\[-\lblskip]
\bi
  \ite Feedbackbögen auswerten: drei Punkte, die Sie für die Seminararbeit
    übernehmen
  \ite \textbf{Folienkritik} abschließen (Aufgabe~4), falls in der Stunde
    nicht fertig
  \ite \textbf{Tagebuch des Staunens}: Wann haben Sie zuletzt jemandem
    geglaubt, weil er überzeugend \emph{wirkte}?
  \ite für LE~13: Was in Ihrem Semesterprojekt würden Sie \emph{nicht}
    veröffentlichen -- und warum?
\ei
\os{\newpage}
\end{description}

\handoutinput{handouts/LE12}
